\makeabstract
    {
    Bird Feeders are Bear Feeders: Studying the Behavior of Ursus americanus in Massachusetts Using Citizen Science
    }
    {
    Faye Riebe\textsuperscript{1}, Leilani Fahn\textsuperscript{1}, Aleksei Sample-Kietrys\textsuperscript{1}, Thea Kristensen\textsuperscript{1}
    }
    {
    \textsuperscript{1} Department of Biology, Amherst College, Amherst, MA
    }
    {
    Increasing urbanization leads to an increase in human-wildlife interaction. One way to monitor wildlife behavior in human-occupied areas is through citizen science, which involves community members submitting data. Here, we sought to explore what behaviors bears (Ursus americanus) were exhibiting and what types of behaviors people tended to report in Massachusetts. Observations may be biased positively, as citizen science is often motivated by interest in conservation. Alternatively, people may be likely to report negative interactions due to fear or negative attitudes towards bears.
    }
    {
    Community volunteers submitted photos of black bears in Massachusetts through the MassBears website. We used photos submitted between April 16, 2019 and July 6, 2026. We reviewed photos to produce a coding scheme, similar to an ethogram, which we then used to define observations and behaviors in each photo. Using R, we assessed the proportion of bears exhibiting each type of behavior. Further, we mapped locations of bears in ArcGISPro. 
    }
    {
    Volunteers submitted 1179 photos. Of these photos, 92.7\% contained bears, 6.1\% had evidence of bears, and 1.2\% had damage caused by bears. For bear photos, 15.3\% were either eating in the photo or the volunteer indicated that they were eating. Of submissions showing the bear eating or with the volunteer indicating that the bear was eating, 52.5\% of submissions indicated the bear was eating bird food. Most sightings of bears were submitted during the summer, whereas the fewest sightings were submitted during the winter. Bears similarly were observed eating bird food most often in summer months. Of the 1,093 bear photos, 117 showed a bear eating.
    }
    {
    The vast majority of sightings were bears, rather than evidence of bears or damage caused by bears. This suggests that fear or negative feelings about bears were not the motivating factor in reporting sightings. It is possible that these types of sightings indicate a more positive attitude towards bears. However, this necessitates further investigation. Gaining a better understanding of the behaviors of black bears in Massachusetts can assist in reducing conflict and fostering human-wildlife coexistence.
    }

    \newpage
    